\specialchapter{参考文献}

\begin{enumerate}
\def\labelenumi{\arabic{enumi}.}
\item
  L. EULER, 代数的完整教程 (=Opera Omnia (1), 第~I 卷，Leipzig-Berlin (Teubner)，1911).
\item
  J. L. LAGRANGE, 著作，14 卷，Paris (Gauthier-Villars)，1867--1892.
\item
  C. F. GAUSS, 著作，12 卷，Göttingen，1870--1927.
\item
  P. G. LEJEUNE-DIRICHLET, 著作，2 卷，Berlin (Reimer)，1889-1897.
\end{enumerate}

4 (bis). P. G. LEJEUNE-DIRICHLET, 数论讲义，第 2 版，Braunschweig (Vieweg)，1871.

\begin{enumerate}
\def\labelenumi{\arabic{enumi}.}
\setcounter{enumi}{4}
\item
  C. G. J. JACOBI, 著作集，7 卷，Berlin (Reimer)，1881-1891.
\item
  G. EISENSTEIN: (a) 关于由单位根的三次根组成的数的理论中三次剩余互反律的证明，Crelle's Journal，27 (1844)，第~289--310 页；(b) 关于素数 8n + 3、7n + 2 和 7n + 4 的二次分解理论，Crelle's Journal，37 (1848)，第~97--126 页；(c) 关于决定整个双纽线分割的方程的一些一般性质及其在数论中的应用，Crelle's Journal，39 (1850)，第~160--179 页和 224--287 页。
\item
  E. KUMMER: (a) 关于由实整数和单位根组成的复数，J. de Math.，(1)，12 (1847)，第~185--212 页；(b) 关于复数理论，Crelle's Journal，35 (1847)，第~319--326 页；(c) 关于由单位根构成的复数分解为素因子，Crelle's Journal，35 (1847)，第~327--367 页；(d) 关于由单位根和整数构成的复数的论文，J. de Math.，(1)，16 (1851)，第~377--498 页；(e) 关于次数为素数的幂的剩余与非剩余之间的一般互反律 (Abh. der Kön. Akad. der Wiss. zu Berlin (1859)，Math. Abhandl.，第~19--159 页)。
\item
  C. HERMITE, 著作，4 卷，Paris (Gauthier-Villars)，1905--1917.
\item
  L. KRONECKER, 著作，5 卷，Leipzig (Teubner)，1895--1930: (a) 关于复单位，第~I 卷，第~5--71 页 (=Inaug. Diss., Berolini，1845); (b) 关于代数可解方程 I，第~IV 卷，第~1--11 页 (=Monatsber. der Kön. Preuss. Akad. der Wiss.，1853，第~365--374 页); (c) 关于具有复乘法的椭圆函数，第~IV 卷，第~177--183 页 (=Monatsber. der Kön. Preuss. Akad. der Wiss.，1857，第~455--460 页); (d) 关于椭圆函数的复乘法，第~IV 卷，第~207--217 页 (=Monatsber. der Kön. Preuss. Akad. der Wiss.，1862，第~363--372 页); (e) 关于一个变量的代数函数的判别式，第~II 卷，第~193--236 页 (=Crelle's Journal，91 (1881)，第~301--334 页); (f) 代数数量的
\end{enumerate}

算术理论概要，第~II 卷，第~237--387 页 (=Crelle's Journal，92 (1882)，第~1--122 页)。

\begin{enumerate}
\def\labelenumi{\arabic{enumi}.}
\setcounter{enumi}{9}
\tightlist
\item
  R. DEDEKIND, 数学著作集，3 卷，Braunschweig (Vieweg)，1932: (a) 关于实素数模的高次同余理论概要，第~I 卷，第~40--66 页 (=Crelle's Journal，54 (1857)，第~1--26 页；(b) 关于代数整数理论，第~III 卷，第~262--296 页 (=Bull. Sci. Math.，(1)，11 (1876)，第~278--288 页和 (2)，1 (1877)，第~17--41 页、69--92 页、144--164 页、207--248 页); (c) 关于有限域不同阶中的理想类数，第~I 卷，第~105--157 页 (=Festschrift der Technischen Hochschule in Braunschweig zur Säkularfeier des Geburtstages von C. F. Gauss，Braunschweig，1877，第~1--55 页); (d) 关于理想理论与高次同余理论之间的联系，第~I 卷，第~202--230 页 (=Abh. Kön. Ges. Wiss. zu Göttingen，23 (1878)，第~1--23 页); (e) 关于有限域的判别式，第~I 卷，第~351--396 页 (=Abh. Kön. Ges. Wiss. zu Göttingen，29 (1882)，第~1--56 页); (f) 关于整代数数理论，第~III 卷，第~1--222 页 (=Supplement XI von Dirichlets Vorlesungen über Zahlentheorie，4 Aufl. (1894)，第~434--657 页); (g) 关于理想理论，第~II 卷，第~43--48 页 (=Nachr. Göttingen，1894，第~272--277 页); (h) 关于模理论中符号 (a, b) 的扩展，第~II 卷，第~59--85 页 (=Nachr. Göttingen，1895，第~183--208 页)。
\end{enumerate}

10 (bis). R. DEDEKIND-H. WEBER, 一个变量的代数函数理论，Crelle's Journal，92 (1882)，第~181--290 页 (=R. Dedekind, Ges. Math. Werke，第~I 卷，第~238--349 页)。

\begin{enumerate}
\def\labelenumi{\arabic{enumi}.}
\setcounter{enumi}{10}
\item
  E. SELLING, 关于由任意不可约方程的根有理构成的复数的理想素因子，Zeitschr. für Math. und Phys.，10 (1865)，第~17--47 页。
\item
  M. NOETHER, 关于代数函数理论中的一个定理，Math. Ann.，6 (1873)，第~351--359 页。
\item
  A. BRILL-M. NOETHER, 关于代数函数，Math. Ann.，7 (1874)，第~269--310 页。
\item
  G. ZOLOTAREFF, 关于复数理论，J. de Math. (3)，6 (1880)，第~51--84 页和 129--166 页。
\item
  E. NETTO, 关于消元理论，Acta Math.，7 (1885)，第~101-104 页。
\item
  D. HILBERT: (a) 关于代数型理论，Math. Ann.，36 (1890)，第~473--534 页；(b) 关于完备不变式系，Math. Ann.，42 (1893)，第~313--373 页；(c) Galois 数域理论概要，Gött. Nachr.，(1894)，第~224--236 页；(d) 数论报告，Jahresber. der D. M. V.，4 (1897)，第~175--546 页（由 A. Lévy 和 Th. Got 译成法文，题为 ``代数数域理论''，Paris (Hermann)，1913）。
\item
  K. WEIERSTRASS, 数学著作，7 卷，Berlin (Mayer und Müller)，1894--1927。
\item
  K. HENSEL: (a) 关于代数数理论的新基础，Jahresber. der D. M. V.，6 (1899)，第~83--88 页；(b) 关于代数域的基本方程与非本质判别因子，Gött. Nachr.，(1897)，第~254--260 页；(c) 算术的新基础，Crelle's Journal，127 (1902)，第~51--84 页；(d) 关于代数数和超越数的算术性质，Jahresber. der D. M. V.，14 (1905)，第~545--558 页；(e) 关于数的算术性质，Jahresber. der D. M. V.，16 (1907)，第~299--319 页、388--393 页、474--496 页；(f) 代数数理论，Leipzig (Teubner)，1908。
\item
  E. LASKER, 关于模与理想理论，Math. Ann.，60 (1905)，第~20--116 页。
\item
  E. STEINITZ: (a) 域的代数理论，Crelle's Journal，137 (1910)，第~167--308 页；(b) 代数数域中的矩形系统与模，Math. Ann. 71 (1912)，第~328--354 页和 72 (1912)，第~297--345 页。
\item
  F. S. MACAULAY, 关于给定模系统分解为准素系统，包括 Hilbert 数的某些性质，Math. Ann.，74 (1913)，第~66--121 页。
\item
  J. KÜRSCHAK, 关于极限构造与一般域理论，Crelle's Journal，142 (1913)，第~211--253 页。
\item
  A. FRAENKEL, 关于零因子与环的分解，Crelle's Journal，145 (1914)，第~139--176 页。
\item
  A. OSTROWSKI, 关于函数方程 \(\phi(x)\phi(y) = \phi(x.y)\) 的一些解，Acta Math.，41 (1917)，第~271-284 页。
\item
  E. NOETHER: (a) 环域中的理想理论，Math. Ann.，83 (1921)，第~24--66 页；(b) 代数数域和函数域中理想理论的抽象构造，Math. Ann.，96 (1927)，第~26--61 页。
\item
  H. HASSE: (a) 关于有理数域中用二次型表示数，Crelle's Journal，152 (1923)，第~129--148 页；(b) 关于有理数域中二次型的等价性，Crelle's Journal，152 (1923)，第~205--224 页；(c) Kurt Hensel 对发现局部--整体原理的决定性推动，Crelle's Journal，209 (1960)，第~3--4 页。
\item
  H. JUNG, 代数曲面，Hannover (Helwing)，1925。
\item
  H. GRELL, 不同环的理想之间的关系，Math. Ann.，97 (1927)，第~490-523 页。
\item
  W. KRULL: (a) 一般环域中的素理想链，Sitz. Ber. Heidelberg Akad. Wiss.，1928；(b) 一般赋值理论，Crelle's Journal，167 (1931)，第~160--196 页；(c) 交换整环算术的贡献，III，Math. Zeitschr.，42 (1937)，第~745--766 页；(d) 位环中的维数理论，Crelle's Journal，179 (1938)，第~204--226 页；(e) 理想理论，Berlin (Springer)，1935。
\item
  M. H. STONE: (a) 布尔代数的表示理论，Trans. Amer. Math. Soc.，40 (1936)，第~37-111 页；(b) 布尔环理论在一般拓扑中的应用，Trans. Amer. Math. Soc.，41 (1937)，第~375-481 页。
\item
  B. L. van der WAERDEN, 现代代数，第~II 卷，Berlin (Springer)，1931。
\item
  C. CHEVALLEY: (a) 无穷扩张的类域理论推广，J. de Math.，(9)，15 (1936)，第~359--371 页；(b) 局部环理论，Ann. of Math.，44 (1943)，第~690--708 页。
\item
  O. ZARISKI: (a) 抽象代数函数域的黎曼流形的紧性，Bull. Amer. Math. Soc.，50 (1944)，第~683--691 页；(b) 广义半局部环，Summa Bras. Math.，1 (1946)，第~169--195 页。
\item
  N. JACOBSON, 任意环的本原理想集上的拓扑，Proc. Nat. Acad. Sci. U.S.A.，31 (1945)，第~333--338 页。
\item
  I. COHEN-A.SEIDENBERG, 素理想与整独立性，Bull. Amer. Math. Soc.，52 (1946)，第~252--261 页。
\item
  P. SAMUEL, 代数与代数几何中的重数概念，J. de Math.，(9)，30 (1951)，第~159--274 页。
\item
  A. WEIL, 代数几何中的纤维空间（A. Wallace 笔记），Chicago Univ.，1952。
\item
  J. P. SERRE: (a) 凝聚代数层，Ann. of Math.，61 (1955)，第~197--278 页；(b) 代数几何与解析几何，Ann. Inst. Fourier，6 (1956)，第~1--42 页。
\item
  A. GROTHENDIECK, 代数几何原理，Publ. math. Inst. Htes Et. Scient.，1960。
\end{enumerate}

